\section{战争不是三条战线，而是三套补给问题}
\label{sec:warfare-diplomacy}

\subsection{夷陵以后，联盟是一种节约}

二二一年刘备东下的军队把失荆州、关羽之死和蜀汉集团的压力带到江汉；
孙吴不能只等魏的册命来决定自己的位置。二二二年，陆逊在夷陵、猇亭击破
蜀军，\eventanchor{TK-0222-YILING}{夷陵之战}改变了三方的选择空间。战役
的结果和刘备退向白帝可由基本叙事相互校验；火攻的过程、营垒延展与伤亡数
却没有同样牢靠的证据。较稳的解释是，蜀军深入后补给线拉长，吴军依水陆
节点等待对手失去机动；它不需要借戏剧化的细节才能成立。

二二三年邓芝使吴，\eventanchor{TK-0223-SHU-WU-RECONCILIATION}{孙刘复通}。
这份关系并未消除荆州的裂痕，而是一种成本计算：蜀若继续在南中、汉中和
江东三方向承压，吴若单独面对北方仓储和陆军，双方都可能先让魏获得决定性
优势。联盟使蜀汉能够转向南中和汉中，使吴把兵力留在江面；它没有创造一支
统一军队，也随时受使节、地方将领与边境城戍的限制。

\subsection{山道、辽东与海面不是背景}

二三〇至二三一年曹真自关中准备攻蜀，栈道和山道遇大雨而受阻；
\eventanchor{TK-0230-CAO-ZHEN-HANZHONG}{关中攻蜀受阻}提醒人们，魏可以
集中人力，却不能把关中仓储直接穿过秦岭。蜀军同样在祁山和五丈原受粮道
限制。山路不是战争地图上的空白，而是决定一项军令能否变成连续作战的财政
与劳役问题：谁修栈道、谁运粮、谁在季节变化前返还，都要由地方社会承担。

东北的选择也由距离限定。二三二年公孙渊通孙吴，次年杀吴使而向魏请和；
孙吴的海上册封不能替代陆上援军。二三八年司马懿抵襄平，便能以近距离军力
结束这种平衡。倭女王卑弥呼遣使至魏，则显示册封、礼物和海路信息可以越过
政权边界，但不应被画成魏的直接统治范围。二四四至二四五年毌丘俭攻高句丽
并破丸都城，短期深入山地同样不等于长期行政占有。

\subsection{淮南是三方都无法忽略的接口}

淮河方向把魏的近畿、吴的水军与地方军镇接在一起。二四四年曹爽攻兴势而退，
二四七年姜维借陇西羌胡冲突出兵，说明西部和南部前线都要在山地、守险和
地方关系中计算。二五二年孙吴在东兴守住水陆节点，二五三年合肥新城远征
又因疫病与补给撤回；同一条水网可以降低防守成本，却不保证进攻者的粮道。

魏内部的淮南叛乱进一步把外战与中枢政治连在一起。二五五年毌丘俭、文钦
起兵，二五七至二五八年诸葛诞据寿春并向吴求援；
\eventanchor{TK-0258-ZHUGE-DAN-DEFEAT}{寿春陷落}后，司马昭压住吴援，也
取得镇压外镇的信誉。三次行动的政治诉求和人际网络不同，不能写成一场连贯
的“复魏”战争；共同点是掌兵者都能利用淮南的道路、城池和跨境联络，迫使
洛阳把中军、转输和人事控制握得更紧。

\begin{turningpoint}[战争的反馈回路]
三国长期均势不是任何一方都没有力量，而是各方把力量运到对手面前的成本都
很高。夷陵促成孙刘复通，复通使魏须维持西、南两线；山道和江面抬高远征
成本，继承与辅政又让一次撤军转化为朝廷清洗。二六三年魏能集中多路伐蜀，
并非突然发现一条道路，而是此前的中军整合、关中经营和蜀汉资源压力终于在
同一时点叠合。
\end{turningpoint}

\begin{sourcenote}
战争年序以\sourcebook{三国志}及\sourcebook{资治通鉴}为主；辽东、高句丽
和倭的部分须与区域史、山城考古及地理研究并读。兵额、斩获、死亡、精确路线
和“完全控制”的说法，在没有交叉证据时仅能作为史书声明，而非本卷的事实
尺度。四幅图只标出阅读所需的相对走廊，不替代历史地理复原。
\end{sourcenote}
